\section*{第 8 讲\quad Tate 上同调与逆向极限}

回忆：对局部域的扩张 $L/K$、Galois 群 $G:=\Z/n\Z$，我们试图证明 $\#\HH^0(G,L^\times)=n$。我们曾断言 $\chi(L^\times)=n$ 当且仅当 $\chi(\OK_L^\times)=1$，其中 $\chi$ 表示 Herbrand 商 $\#\HH^0/\#\HH^1$；并回忆有限群的 Herbrand 商等于 $1$，且 $\chi$ 对短正合列具有乘性。

上一次我们证明了 $\chi(\OK_L)=1$：用正规基定理说明 $L\simeq K[G]$、$\OK_L$ 含有一个有限指数的开子群 $\Gamma$ 使得 $\OK_L^\times\supset\Gamma\simeq\OK_L^+$ 是 $G$-等价的（从而 $\Gamma$ 在 $G$-作用下封闭），并且 $\OK_K[G]\simeq\Gamma\subseteq\OK_L$。随后我们用断言 7.10 说明对每个 $i$ 都有 $\HH^i(G,\Gamma)=0$，从而 $\chi(\OK_L)=\chi(\Gamma)=1$。

现在我们要给出更好、即更代数化（不用 $p$ 进指数映射！）的 $\chi(\OK_L^\times)=1$ 的证明。为此固定一个同构于 $\OK_K[G]$（作为 $\OK_K[G]$-模）的开子群 $\Gamma\subseteq\OK_L$。

\textbf{断言 8.1.} 对充分大的 $N$，$1+\mathfrak p_K^N\Gamma$ 是 $\OK_L^\times$ 的子群，其中 $\mathfrak p_K$ 是 $\OK_K$ 的极大理想。

\textit{证明.} 对 $x,y\in\Gamma\subseteq\OK_L$，有
\begin{equation}\tag{8.1}
(1+\pi^Nx)(1+\pi^Ny)=1+\underbrace{\pi^N(x+y)}_{\in\,\mathfrak p_K^N\Gamma}+\underbrace{\pi^{2N}(xy)}_{\in\,\mathfrak p_K^{2N}\OK_L},
\end{equation}
其中 $\pi$ 是 $\mathfrak p_K$ 的素元。于是，若取 $N$ 充分大使得 $\mathfrak p_K^N\OK_L\subseteq\Gamma$（这是可能的，因为 $\Gamma$ 是 $\OK_L$ 的开子群），则该乘积落在 $1+\pi^N\Gamma$ 中，从而 $1+\mathfrak p_K^N\Gamma$ 是 $\OK_L^\times$ 的子群。\hfill$\square$

\textbf{断言 8.2.} $\Gamma$ 的各阶上同调全部消没。

取 $N$ 使得 $\mathfrak p_K^N\OK_L\subseteq\mathfrak p_K\Gamma$。则 (8.1) 中最后一项落在 $\mathfrak p_K^{2N}\OK_L\subseteq\mathfrak p_K^{N+1}\Gamma$ 中。这表明我们应当用加法商群来过滤 $1+\mathfrak p_K^N\Gamma$，即用 $1+\mathfrak p_K^{N+i}\Gamma$ 过滤，使得由上述计算，对一切 $i\ge 0$ 有加法群的同构
\[
(1+\mathfrak p_K^{N+i}\Gamma)/(1+\mathfrak p_K^{N+i+1}\Gamma)\simeq\Gamma/\mathfrak p_K\Gamma\simeq k_K[G],
\]
其中 $k_K$ 表示 $K$ 的剩余域。此外，这些同构是 Galois 等变的，即 $G$-等变的：$G$-作用保持所有项（$\Gamma$ 由假设保持），从而作用于两边，且被同构保持。于是由断言 7.10，
\[
\HH^j\big(G,(1+\mathfrak p_K^{N+i}\Gamma)/(1+\mathfrak p_K^{N+i+1}\Gamma)\big)=\HH^j(G,k_K[G])=0
\]
对每个 $i\ge 0$ 与 $j$ 成立。作为推论（为此需要下面的引理），
\begin{equation}\tag{8.2}
\HH^j\big(G,(1+\mathfrak p_K^N\Gamma)/(1+\mathfrak p_K^{N+i}\Gamma)\big)=0
\end{equation}
对一切 $i\ge 0$ 与 $j$ 成立。

\textbf{引理 8.3.} 对任意 $G$-模的短正合列
\[
0\to M\to E\to N\to 0,
\]
若 $\HH^i(G,M)=\HH^i(G,N)=0$，则对每个 $i$ 有 $\HH^i(G,E)=0$。

\textit{证明.} 由 (6.4) 有正合列
\[
\HH^i(G,M)\xrightarrow{\alpha}\HH^i(G,E)\xrightarrow{\beta}\HH^i(G,N),
\]
故 $\HH^i(E)=\ker(\beta)=\im(\alpha)=0$，如所欲证。\hfill$\square$

现在有正合列
\[
0\to\frac{1+\mathfrak p_K^{N+i}\Gamma}{1+\mathfrak p_K^{N+i+1}\Gamma}\to\frac{1+\mathfrak p_K^{N}\Gamma}{1+\mathfrak p_K^{N+i+1}\Gamma}\to\frac{1+\mathfrak p_K^{N}\Gamma}{1+\mathfrak p_K^{N+i}\Gamma}\to 0,
\]
于是 (8.2) 可由对 $i$ 的归纳与引理 8.3 得到。

剩下的问题是证明 $\HH^j(1+\mathfrak p_K^N\Gamma)=0$。在理想的情形下，对任意带 $G$-作用与 $G$-等变结构映射的模序列，我们都应有
\[
\HH^j\Big(G,\varprojlim_n M_n\Big)=\varprojlim_n\HH^j(G,M_n).
\]
这样便可推出
\[
\HH^j(1+\mathfrak p_K^N\Gamma)
=\HH^j\Big(\varprojlim_{i\ge 0}(1+\mathfrak p_K^N\Gamma)/(1+\mathfrak p_K^{N+i}\Gamma)\Big)
=\varprojlim_{i\ge 0}\HH^j\big((1+\mathfrak p_K^N\Gamma)/(1+\mathfrak p_K^{N+i}\Gamma)\big)
=\varprojlim_{i\ge 0}0=0
\]
（由 (8.2)，因为我们的过滤是完备的）。因此，我们需要设法说明 Tate 上同调与逆向极限可以交换。

\textbf{引理 8.4.} 设有一列模
\begin{equation}\tag{8.3}
\begin{CD}
\vdots @. \vdots @. \vdots @. @.\\
0 @>>> M_{n+1} @>>> E_{n+1} @>>> N_{n+1} @>>> 0\\
@. @VVV @VVV @VVV @.\\
0 @>>> M_{n} @>>> E_{n} @>>> N_{n} @>>> 0\\
\vdots @. \vdots @. \vdots @. @.
\end{CD}
\end{equation}
其各行正合。则我们有正合列
\[
0\to\varprojlim_n M_n\xrightarrow{\varphi}\varprojlim_n E_n\xrightarrow{\psi}\varprojlim_n N_n .
\]
% [原文疑误：原文作 "if $M_n\to M_{n+1}$ is surjective for each $n$"，与 (8.3) 阶梯图的方向相反；译文按图示取 $M_{n+1}\to M_n$]
此外，若对每个 $n$，$M_{n+1}\to M_n$ 都是满射，则 $\psi$ 也是满射（否则可能不是！）。

\textit{证明.} 显然 $\varphi$ 是单射：若 $x\in\ker(\varphi)$，则其像的每个坐标都是 $0$，故由每个 $M_n\to E_n$ 的单射性与相容性，$x$ 的每个坐标都是 $0$，从而 $x=0$。同理，由 (8.3) 每行正合可得 $\ker(\psi)=\im(\varphi)$。

为（直观地）看出 $\psi$ 的满射性何以会失败，考虑相容系统 $(x_n)\in\varprojlim_n N_n$。我们可以把每个 $x_n$ 提升为某个 $y_n\in E_n$，但如何相容地做到这一点、使 $(y_n)\in\varprojlim_n E_n$，并不清楚。

现在假设每个映射 $M_{n+1}\to M_n$ 都是满射。取 $(x_n)\in\varprojlim_n N_n$，并设对某个 $n$ 已构造出 $y_n\in E_n$。任取 $\tilde y_{n+1}\mapsto x_{n+1}$（经 $E_{n+1}\to N_{n+1}$），令 $\alpha_{n+1}$ 为 $\tilde y_{n+1}$ 在 $E_n$ 中的像。则 $y_n-\alpha_{n+1}\in M_n$（因为它在 $N_n$ 中消没），且由假设它可提升为 $\beta_{n+1}\in M_{n+1}$。若令 $y_{n+1}:=\tilde y_{n+1}+\beta_{n+1}$，则它在 $E_n$ 中映到 $\alpha_{n+1}+y_n-\alpha_{n+1}=y_n$，在 $N_{n+1}$ 中映到 $x_{n+1}$（因为 $\beta_{n+1}$ 在 $N_{n+1}$ 中映到 $0$）；于是由归纳法，存在相容系统 $(y_n)\in\varprojlim_n E_n$ 经 $\psi$ 映到 $(x_n)$，如所欲证（注意，这一结果可表述为满射 $E_{n+1}\twoheadrightarrow E_n\times_{N_n}N_{n+1}$，即到纤维积）。\hfill$\square$

\textbf{命题 8.5.} 若
\begin{equation}\tag{8.4}
\cdots\twoheadrightarrow M_{n+1}\twoheadrightarrow M_n\twoheadrightarrow\cdots\twoheadrightarrow M_0
\end{equation}
是 $G$-模的序列（各映射均为满射；证明中要用到这一点，见下文），且对一切 $n$ 与 $i$ 有 $\HH^i(M_{n+1})\twoheadrightarrow\HH^i(M_n)$，则对一切 $i$ 有
\[
\HH^i\Big(\varprojlim_n M_n\Big)=\varprojlim_n\HH^i(M_n).
\]

\textit{证明.} 我们对 $\HH^0$ 给出证明。令 $M:=\varprojlim_n M_n$，于是我们要比较 $\HH^0(M)=M^G/N(M)$ 与 $\varprojlim_n\HH^0(M_n)=\varprojlim_n M_n^G/N(M_n)$。这归结为证明自然映射
\[
\Big(\varprojlim_n M_n^G\Big)\Big/\Big(\varprojlim_n N(M_n)\Big)\xrightarrow{\sim}\varprojlim_n\big(M_n^G/N(M_n)\big)
\]
是同构。我们有交换图
\[
\begin{CD}
0 @>>> N(M_{n+1}) @>>> M_{n+1}^G @>>> \HH^0(G,M_{n+1}) @>>> 0\\
@. @V{\alpha_n}VV @VVV @VVV @.\\
0 @>>> N(M_{n}) @>>> M_{n}^G @>>> \HH^0(G,M_{n}) @>>> 0,
\end{CD}
\]
并断言每个 $\alpha_n$ 都是满射。事实上，取 $x\in N(M_n)$，则对某个 $y\in M_n$ 有 $x=N(y)$。把 $y$ 提升为元素 $z\in M_{n+1}$，便有 $\alpha_n(N(z))=x$，如所欲证。于是由引理 8.4，
\[
\varprojlim_n\HH^0(M_n)=\Big(\varprojlim_n M_n^G\Big)\Big/\Big(\varprojlim_n N(M_n)\Big).
\]
现在
\[
\HH^0\Big(\varprojlim_n M_n\Big)=\Big(\varprojlim_n M_n\Big)^G\Big/N\Big(\varprojlim_n M_n\Big)=M^G/N(M).
\]
由于 $G$ 作用于 $\varprojlim_n M_n$ 的每个坐标，显然 $(\varprojlim_n M_n)^G=\varprojlim_n M_n^G$；于是只剩下证明 $N(\varprojlim_n M_n)\xrightarrow{\sim}\varprojlim_n N(M_n)$。对每个 $n$ 令 $K_n:=\ker(N:M_n\to M_n)$，则有交换图
\[
\begin{CD}
0 @>>> K_{n+1} @>>> M_{n+1} @>>> N(M_{n+1}) @>>> 0\\
@. @VVV @VVV @V{\beta_n}VV @.\\
0 @>>> K_{n} @>>> M_{n} @>>> N(M_{n}) @>>> 0.
\end{CD}
\]
我们希望证明 $\beta_n$ 是满射。回忆 $\HH^1(G,M_n)=K_n/(1-\sigma)M_n$，于是有交换图
\[
\begin{CD}
0 @>>> (1-\sigma)M_{n+1} @>>> K_{n+1} @>>> \HH^1(G,M_{n+1}) @>>> 0\\
@. @V{\gamma_n}VV @V{\beta_n}VV @V{\delta_n}VV @.\\
0 @>>> (1-\sigma)M_{n} @>>> K_{n} @>>> \HH^1(G,M_{n}) @>>> 0.
\end{CD}
\]
现在 $\gamma_n$ 是满射（证明与 $\alpha_n$ 的情形类似），而 $\delta_n$ 由假设是满射。于是由蛇引理，$\beta_n$ 是满射，从而引理 8.4 给出
\[
\varprojlim_n N(M_n)=\Big(\varprojlim_n M_n\Big)\Big/\Big(\varprojlim_n K_n\Big)=N\Big(\varprojlim_n M_n\Big).
\]
由此即得 $\HH^0(\varprojlim_n M_n)=\varprojlim_n\HH^0(M_n)$，如所欲证。\hfill$\square$

\textbf{推论 8.6.} 对序列 (8.4)，若对一切 $n$ 与 $i$ 有 $\HH^i(M_n)=0$，则
\[
\HH^i\Big(\varprojlim_n M_n\Big)=0 .
\]
特别地，$\HH^i(1+\mathfrak p_K^N\Gamma)=0$，其中我们置
\[
M_i:=(1+\mathfrak p_K^N\Gamma)/(1+\mathfrak p_K^{N+i}\Gamma)\qquad(i\ge 0).
\]
于是，由于 $1+\mathfrak p_K^N\Gamma\subseteq\OK_L^\times$ 是有限指数子群（作为加法群带有正规基），我们有
\[
\chi(\OK_L^\times)=\chi(1+\mathfrak p_K^N\Gamma)=1,
\]
这就建立了断言 7.8 的 (2)。
